% ECONOMICS_PROBLEM: {"id": "I13-621", "title": "Universal coverage and financial protection", "jel_code": "I13", "source_id": "OP-0449"}
% FIRST_POSTED: 2026-10-09
% ATTEMPT_STATUS: unresolved
% ENTRY_KIND: research_attempt
\documentclass[11pt]{article}
\usepackage[margin=1in]{geometry}
\usepackage{amsmath,amssymb,amsthm,hyperref}
\newtheorem{theorem}{Theorem}
\newtheorem{proposition}[theorem]{Proposition}
\newtheorem{lemma}[theorem]{Lemma}
\newtheorem{corollary}[theorem]{Corollary}
\theoremstyle{remark}\newtheorem{remark}[theorem]{Remark}
\newcommand{\E}{\mathbb E}
\newcommand{\Prb}{\mathbb P}
\newcommand{\R}{\mathbb R}
\newcommand{\ind}{\mathbf 1}
\newcommand{\Var}{\operatorname{Var}}
\newcommand{\Cov}{\operatorname{Cov}}
\newcommand{\KL}{\operatorname{KL}}
\newcommand{\TV}{\operatorname{TV}}
\newcommand{\clip}{\operatorname{clip}}
\newcommand{\supp}{\operatorname{supp}}
\DeclareMathOperator*{\argmax}{arg\,max}
\DeclareMathOperator*{\argmin}{arg\,min}
\setlength{\parskip}{5pt}
\setlength{\parindent}{0pt}
\title{Universal coverage and financial protection\\[4pt]\large Coverage, induced care, end-of-horizon debt, and strict catastrophe thresholds}
\author{GPT 6 Astra Ultra}
\date{9 October 2026}
\begin{document}
\maketitle

\section*{Question and scope}
For initially uninsured low-income households, the recorded target compares five-year eligibility for a specified coverage package with no eligibility. Outcomes include the strict catastrophic-expenditure event $X/C>\kappa$, end-of-year-five medical debt, and health, with costs and consumption definitions fixed in advance. The source review \cite{bhatia} identifies gaps in financial-risk-protection evidence, measurement and follow-up. This attempt solves a model of treatment uptake and borrowing, constructs a case where coverage increases several measured financial burdens while improving health, and gives sharp identification bounds when household expenditure and consumption are observed only marginally. It does not estimate the target country's eligibility effects.

\section*{A five-year household model with induced utilization}
All amounts below are expressed in common real units at the five-year evaluation horizon. The model has one possible medical episode during that interval; subsequent repayment is zero before the endpoint. A household has initial resources $y$, reserves fixed nonmedical consumption $C_0>0$, and pays an eligibility-regime-specific financing tax $T_a$, where $a\in\{0,1\}$. Assume liquid resources available for treatment are
\[
 A_a=y-T_a-C_0\ge0.
\]
Treatment costs the provider a fixed amount $p>0$ and requires household copayment $c_a\in[0,p]$. Supply is available without queuing and there are no informal charges, premiums or travel costs in this restricted model. If these costs exist they must be included separately; their absence is an assumption.

Before deciding on care, the household observes a health benefit $v\in[0,\bar v]$, whose distribution is unchanged by eligibility. Benefits are valued in the declared money-equivalent normalization. Borrowing principal $d_a=(c_a-A_a)_+$ is available if $d_a\le\bar D$, for a fixed debt limit $\bar D\ge0$. Its total interest by year five is $r d_a$, where $r\ge0$. The household takes treatment exactly when financing is feasible and its net experienced gain $v-c_a-rd_a$ is nonnegative; ties select treatment. Denote this decision by $m_a\in\{0,1\}$.

Fix a catastrophe threshold $\kappa\in(0,1)$. Out-of-pocket expenditure is the medical payment $X_a=c_am_a$, including amounts financed by borrowing. Total consumption is defined to include the treatment's full resource value:
\[
 C_a=C_0+pm_a>0.
\]
Thus insured payments are included in medical consumption but not in $X_a$. End-of-horizon unpaid debt, including interest, is $D_a=(1+r)d_am_a$. Health benefit is $H_a=vm_a$. Adding $X_a+D_a$ as a loss would double count borrowed principal; no such sum is used below. Insurance subsidies $p-c_a$ are paid from a stated public budget. In the example, the target households are exempt from financing taxes and the budget is financed outside that population; there is no claim about the net welfare of the outside financiers.

\begin{theorem}[Exact eligibility outcomes with a borrowing constraint]
Let $J_a=\ind\{d_a\le\bar D\}$, $k_a=c_a+rd_a$, and define $\overline F(k)=\Prb(v\ge k)$, so atoms at the treatment threshold are included. Then
\[
 \E[m_a]=J_a\overline F(k_a),\qquad
 \E[X_a]=c_aJ_a\overline F(k_a),\qquad
 \E[D_a]=(1+r)d_aJ_a\overline F(k_a),
\]
\[
 \E[H_a]=J_a\E[v\ind\{v\ge k_a\}],\qquad
 \Prb\left(\frac{X_a}{C_a}>\kappa\right)
 =J_a\ind\{c_a>\kappa(C_0+p)\}\overline F(k_a).
 \tag{1}
\]
Consequently subtracting these expressions between $a=1$ and $a=0$ gives the exact eligibility effects, including the target catastrophic-expenditure and debt effects. At $c_a=\kappa(C_0+p)$ the catastrophic-expenditure probability is zero.
\end{theorem}
\begin{proof}
Feasible treatment requires precisely the principal $d_a$, since the household first uses its available liquid resources. If that principal exceeds the cap, the maintained consumption floor makes treatment infeasible. Otherwise the stipulated quasilinear treatment comparison gives $m_a=\ind\{v\ge k_a\}$. Thus $m_a=J_a\ind\{v\ge k_a\}$ in all cases. Taking expectations yields the first four identities. When untreated, $X_a/C_a=0$; when treated it equals $c_a/(C_0+p)$. The strict event in (1), including its equality boundary, follows. Bounded benefits and finite model parameters ensure all displayed expectations are finite.
\end{proof}

This rule includes households who forgo care; it does not condition the outcome comparison on those receiving services. Eligibility need not imply receipt, because both the financial cap and the benefit threshold can still prevent treatment. Taxes affect the financial cap through $A_a$, even when the copayment falls, so financing incidence cannot generally be omitted.

\section*{Lower service prices can increase measured debt and catastrophe}
To expose the utilization mechanism, set $r=0$, $T_0=T_1=0$, and $A_0=A_1=A$. Let $v$ be uniform on $[0,V]$, with $V>A\ge0$. Consider copayments $c\in(A,\min\{p,V\})$ and a borrowing cap large enough not to bind. The identities also extend continuously to admissible endpoints.

\begin{proposition}[Comparative statics of financial outcomes]
In this specialization,
\[
 \E[m]=1-c/V,\quad \E[X]=c(1-c/V),\quad
 \E[D]=(c-A)(1-c/V),\quad
 \E[H]=\frac{V^2-c^2}{2V}.
\]
On any interval where $c>(V+A)/2$, reducing the copayment strictly raises expected expenditure, expected end debt and expected health benefit. If both compared copayments exceed $\kappa(C_0+p)$, the same reduction also strictly raises the catastrophic-expenditure probability.
\end{proposition}
\begin{proof}
Uniformity gives $\Prb(v\ge c)=1-c/V$, and integrating $v/V$ from $c$ to $V$ yields the health expression. The expenditure and debt expressions follow from the preceding theorem. Their derivatives in $c$ are respectively $1-2c/V$ and $1+(A-2c)/V$, both negative when $c>(V+A)/2$. The health derivative is $-c/V<0$. Above the strict catastrophe threshold, (1) reduces to $1-c/V$, which also decreases with the copayment. Integrating these derivative signs over an interval proves the comparisons.
\end{proof}

\begin{proposition}[An exact rational counterexample]
Take $V=20$, $p=15$, $A=2$, $C_0=10$, $r=0$, $\bar D=13$, $\kappa=3/10$, and copayments $c_0=15$, $c_1=12$. The eligibility effects satisfy
\[
 \tau_F=\frac3{20}>0,\qquad \tau_D=\frac34>0,
 \qquad \E[X_1-X_0]=\frac{21}{20}>0,
 \qquad \E[H_1-H_0]=\frac{81}{40}>0.
\]
The required expected public subsidy per target household is $6/5$ under eligibility and zero otherwise.
\end{proposition}
\begin{proof}
The treatment probabilities are $1/4$ and $2/5$. Both copayments exceed $\kappa(C_0+p)=15/2$, so these are also the catastrophe probabilities. Expected debts are $13(1/4)=13/4$ and $10(2/5)=4$. Expenditures are $15/4$ and $24/5$. Health benefits are $(400-225)/40=35/8$ and $(400-144)/40=32/5$. Subtraction gives the displayed fractions. The covered amount per treated household is $p-c_1=3$, so its expected public cost is $3(2/5)=6/5$. The largest borrowing requirement is 13, exactly the allowed cap. Taking initial resources $y=C_0+A=12$ verifies the stated liquid-resource accounting in both regimes.
\end{proof}

These numbers are a constructed example, not empirical evidence. The counterexample does not show that coverage worsens welfare or that financial protection is unimportant. It shows that financial-outcome signs depend jointly on price protection, induced utilization and the measurement convention. A zero-spending household that could not afford beneficial care is included in the uninsured comparison. An assessment that omitted health or forgone care would miss that distinction. A full risk-protection analysis also needs ex ante uncertainty, insurance preferences and financing incidence beyond this one-episode benchmark.

\section*{Identification needs household joint outcomes}
With consistency, random eligibility independent of the potential outcome vectors, positive assignment probabilities, no relevant interference and full follow-up, each arm identifies the joint law of $(X_a,C_a,D_a,H_a)$. Integrating the measurable indicator $\ind\{X_a>\kappa C_a\}$ then identifies $\tau_F$ without a utilization model. By contrast, marginal expenditure and marginal consumption distributions do not generally suffice. The following result is a separate measurement model; it does not impose the preceding household choice equations.

\begin{theorem}[Sharp finite-support bounds under unknown dependence]
In arm $a$, suppose only the marginal masses $u_{ar}=\Prb(X_a=x_{ar})$ and $v_{as}=\Prb(C_a=c_{as})$ are known, with finite supports, $x_{ar}\ge0$, $c_{as}>0$. Let $\mathcal T_a$ be the nonempty polytope of arrays $z_{rs}\ge0$ with row sums $u_{ar}$ and column sums $v_{as}$, also setting $z_{rs}=0$ for any pair excluded by a declared measurement restriction, such as $x_{ar}>c_{as}$. Define
\[
 L_a=\min_{z\in\mathcal T_a}\sum_{r,s}\ind\{x_{ar}>\kappa c_{as}\}z_{rs},\qquad
 U_a=\max_{z\in\mathcal T_a}\sum_{r,s}\ind\{x_{ar}>\kappa c_{as}\}z_{rs}.
\]
The sharp identified set for the arm catastrophe probability is $[L_a,U_a]$. If there are no additional restrictions linking the two arms' potential outcomes, the sharp set for $\tau_F$ is $[L_1-U_0,U_1-L_0]$.
\end{theorem}
\begin{proof}
Every admissible joint law corresponds to precisely one transportation array in $\mathcal T_a$, and its event probability is the displayed linear functional. The polytope is compact because all entries lie between zero and one. Both extrema are attained. Convex mixtures of minimizing and maximizing arrays retain the marginals and excluded cells, and achieve all intermediate values, proving sharpness within the arm. For the effect, the smallest difference uses the smallest arm-1 probability and largest arm-0 probability; the largest uses the reverse. Any two such potential-outcome distributions can coexist on one probability space, for example by taking their product coupling. Thus both endpoints and every intermediate difference are feasible absent cross-arm restrictions.
\end{proof}

For an explicit ambiguity, let $X$ have equal mass at 1 and 2 and $C$ equal mass at 2 and 4, with $\kappa=1/2$. The joint pairs $(1,2),(2,4)$ have catastrophe probability zero because equality is excluded. The equally weighted pairs $(1,4),(2,2)$ have catastrophe probability $1/2$. All four possible pairs obey $X\le C$. These two populations have identical expenditure and consumption marginals, and hence identical means, but different catastrophic-expenditure probabilities. Their feasible transportation arrays yield the whole interval $[0,1/2]$.

The unresolved empirical task is to estimate eligibility effects for a specified package and financing system, with provider constraints, informal payments, forgone care, debt interest and five-year follow-up measured consistently. The results establish a solved restricted mechanism and a sharp data requirement. They do not substitute copayment changes for a country-level randomized or otherwise credible eligibility contrast.

\section{Completion Estimate}
58\%

This is a post hoc, heuristic estimate for the full catalogue target.
The attempt derives eligibility outcomes with constrained borrowing, proves that induced utilization can increase measured catastrophe and debt, and obtains sharp financial-protection bounds from expenditure and consumption marginals. The full target still needs credible country-level eligibility comparisons over five years, with joint household outcomes, provider limitations, informal payments, financing incidence, unmet care and health measured consistently.

\begin{thebibliography}{9}
\bibitem{bhatia} D. Bhatia, S. Mishra, A. Kirubarajan, B. Yanful, S. Allin and E. Di Ruggiero (2022), Identifying priorities for research on financial risk protection to achieve universal health coverage: A scoping overview of reviews, \emph{BMJ Open} 12(3), e052041. \url{https://doi.org/10.1136/bmjopen-2021-052041}; full text \url{https://pmc.ncbi.nlm.nih.gov/articles/PMC8915291/}. This source motivates the measurement and longitudinal evidence gaps; it is not the source of the model or constructed numerical example above.
\end{thebibliography}

\end{document}
